% GENERATED FILE. Edit canonical lessons, then run npm run generate:books.
% canonical-source-hash: fnv1a64:c735638084118587
% canonical-lessons: LA-C478-opilio, LA-C478-bubulcus, LA-C478-caulae, LA-C478-tondeo, LA-C478-mulgeo, LA-R478-review-471-478

\chapter{Shepherd, Ox-driver, Sheepfold, To shear, To milk}
\label{ch:la-shepherd-ox-driver-sheepfold-to-shear-to-milk}
\hlchaptermodality{478}

\begin{chapteropening}
\textbf{By the end of this chapter:} \emph{I can talk about the work of herdsmen: the shepherd, the ox-driver, the sheepfold, shearing sheep and milking.}

Complete the last lesson of chapter 478: I can talk about the work of herdsmen: the shepherd, the ox-driver, the sheepfold, shearing sheep and milking.
\end{chapteropening}

\section[ōpiliō, ōpiliōnis]{ōpiliō, ōpiliōnis — a shepherd (of sheep)}

\label{lesson:LA-C478-opilio}



\begin{quote}
Before the new one: say the Latin for to pasture, then the Latin for a pasture.
\end{quote}

\subsection*{You'll want to know: ōpiliō, ōpiliōnis}
\textbf{ōpiliō, ōpiliōnis} — \textquotedblleft{}a shepherd (of sheep)\textquotedblright{}.

\begin{quote}
Ōpiliō ovēs in colle servat.
\end{quote}

\begin{quote}
\emph{The shepherd watches over the sheep on the hill.}
\end{quote}

Third declension, masculine: \emph{ōpiliō, ōpiliōnis m.}, a word of comedy and the farm writers. \emph{Pāstor} is the general word.

\begin{cousinweb}
\emph{ōpiliō} is usually traced to \emph{ovis} 'sheep' (Lewis \& Short, from an older \emph{ovilio}), so a sheep-keeper; the second part is not settled. EtymDB lists it without a parent.
\end{cousinweb}

\begin{grammarlens}[title={What you've built}]
One more word for telling, explaining and describing at B1.
\end{grammarlens}

\subsection*{Your turn}
\begin{itemize}
\raggedright
  \item \emph{Say it:} \emph{ōpiliō, ōpiliōnis}
  \item \emph{Say it:} \emph{ōpiliō, ōpiliōnis}, once more
  \item \emph{Say it:} a sentence of your own with \emph{ōpiliō, ōpiliōnis}
  \item \emph{Recall:} say the Latin for to pasture, then the Latin for a pasture, then say \emph{ōpiliō, ōpiliōnis} again
\end{itemize}

\begin{tcolorbox}[breakable,colback=teal!4,colframe=teal!35!black,title={Before you move on}]
What does ōpiliō, ōpiliōnis mean? (\textquotedblleft{}A shepherd (of sheep)\textquotedblright{}.) Say it once more.
\end{tcolorbox}
\section[bubulcus, bubulcī]{bubulcus, bubulcī — an ox-driver, a herdsman}

\label{lesson:LA-C478-bubulcus}



\begin{quote}
Before the new one: say the Latin for a pasture, then the Latin for a shepherd.
\end{quote}

\subsection*{You'll want to know: bubulcus, bubulcī}
\textbf{bubulcus, bubulcī} — \textquotedblleft{}an ox-driver, a herdsman\textquotedblright{}.

\begin{quote}
Bubulcus cum bōbus ad agrum it.
\end{quote}

\begin{quote}
\emph{The ox-driver goes to the field with the oxen.}
\end{quote}

Second declension, masculine: \emph{bubulcus, bubulcī m.} The man who drives the plough oxen.

\begin{cousinweb}
\emph{bubulcus} is built on \emph{būbulus} 'of cattle', from \emph{bōs, bovis} 'ox' (EtymDB; Lewis \& Short), from PIE \emph{g\textsuperscript{w}\'{\={o}}ws} 'cow'. English \emph{cow} is a Germanic cousin.
\end{cousinweb}

\begin{grammarlens}[title={What you've built}]
One more word for telling, explaining and describing at B1.
\end{grammarlens}

\subsection*{Your turn}
\begin{itemize}
\raggedright
  \item \emph{Say it:} \emph{bubulcus, bubulcī}
  \item \emph{Say it:} \emph{bubulcus, bubulcī}, once more
  \item \emph{Say it:} a sentence of your own with \emph{bubulcus, bubulcī}
  \item \emph{Recall:} say the Latin for a pasture, then the Latin for a shepherd, then say \emph{bubulcus, bubulcī} again
\end{itemize}

\begin{tcolorbox}[breakable,colback=teal!4,colframe=teal!35!black,title={Before you move on}]
What does bubulcus, bubulcī mean? (\textquotedblleft{}An ox-driver, a herdsman\textquotedblright{}.) Say it once more.
\end{tcolorbox}
\section[caulae, caulārum]{caulae, caulārum — a sheepfold, a pen; an opening}

\label{lesson:LA-C478-caulae}



\begin{quote}
Before the new one: say the Latin for a shepherd, then the Latin for an ox-driver.
\end{quote}

\subsection*{You'll want to know: caulae, caulārum}
\textbf{caulae, caulārum} — \textquotedblleft{}a sheepfold, a pen; an opening\textquotedblright{}.

\begin{quote}
Lupus prope caulās vagātur.
\end{quote}

\begin{quote}
\emph{A wolf prowls near the sheepfold.}
\end{quote}

First declension, feminine plural: \emph{caulae, caulārum f.} First a hole or passage, then a fenced pen.

\begin{cousinweb}
\emph{caulae} is usually traced to \emph{cavus} 'hollow', through an older \emph{cavile} (Lewis \& Short, after an ancient grammarian); EtymDB gives no parent, so the line is hedged.
\end{cousinweb}

\begin{grammarlens}[title={What you've built}]
One more word for telling, explaining and describing at B1.
\end{grammarlens}

\subsection*{Your turn}
\begin{itemize}
\raggedright
  \item \emph{Say it:} \emph{caulae, caulārum}
  \item \emph{Say it:} \emph{caulae, caulārum}, once more
  \item \emph{Say it:} a sentence of your own with \emph{caulae, caulārum}
  \item \emph{Recall:} say the Latin for a shepherd, then the Latin for an ox-driver, then say \emph{caulae, caulārum} again
\end{itemize}

\begin{tcolorbox}[breakable,colback=teal!4,colframe=teal!35!black,title={Before you move on}]
What does caulae, caulārum mean? (\textquotedblleft{}A sheepfold, a pen; an opening\textquotedblright{}.) Say it once more.
\end{tcolorbox}
\section[tondeō, tondēre]{tondeō, tondēre — to shear, to clip; to shave}

\label{lesson:LA-C478-tondeo}



\begin{quote}
Before the new one: say the Latin for an ox-driver, then the Latin for a sheepfold.
\end{quote}

\subsection*{You'll want to know: tondeō, tondēre}
\textbf{tondeō, tondēre} — \textquotedblleft{}to shear, to clip; to shave\textquotedblright{}.

\begin{quote}
Vēre agricolae ovēs tondent.
\end{quote}

\begin{quote}
\emph{In spring the farmers shear the sheep.}
\end{quote}

\emph{tondeō, tondēre, totondī, tōnsum}, with a reduplicated perfect. \emph{Tōnsor} 'barber' comes from it, and English \emph{tonsure}.

\begin{cousinweb}
\emph{tondeō} is inherited and is usually traced to PIE \emph{temh\textsubscript{2}-} 'to cut' (EtymDB; Lewis \& Short compare Greek \emph{temnō} 'to cut'). English \emph{tonsure} comes from its participle \emph{tōnsus}.
\end{cousinweb}

\begin{grammarlens}[title={What you've built}]
One more word for telling, explaining and describing at B1.
\end{grammarlens}

\subsection*{Your turn}
\begin{itemize}
\raggedright
  \item \emph{Say it:} \emph{tondeō, tondēre}
  \item \emph{Say it:} \emph{tondeō, tondēre}, once more
  \item \emph{Say it:} a sentence of your own with \emph{tondeō, tondēre}
  \item \emph{Recall:} say the Latin for an ox-driver, then the Latin for a sheepfold, then say \emph{tondeō, tondēre} again
\end{itemize}

\begin{tcolorbox}[breakable,colback=teal!4,colframe=teal!35!black,title={Before you move on}]
What does tondeō, tondēre mean? (\textquotedblleft{}To shear, to clip; to shave\textquotedblright{}.) Say it once more.
\end{tcolorbox}
\section[mulgeō, mulgēre]{mulgeō, mulgēre — to milk}

\label{lesson:LA-C478-mulgeo}



\begin{quote}
Before the new one: say the Latin for a sheepfold, then the Latin for to shear.
\end{quote}

\subsection*{You'll want to know: mulgeō, mulgēre}
\textbf{mulgeō, mulgēre} — \textquotedblleft{}to milk\textquotedblright{}.

\begin{quote}
Ancilla mane caprās mulget.
\end{quote}

\begin{quote}
\emph{In the morning the maid milks the goats.}
\end{quote}

\emph{mulgeō, mulgēre, mulsī, mulsum} or \emph{mulctum}. English \emph{milk} is a Germanic cousin, and \emph{emulsion} comes from \emph{ēmulgēre}.

\begin{cousinweb}
\emph{mulgeō} is inherited from PIE \emph{h\textsubscript{2}melǵ-} 'to milk' (EtymDB); German \emph{melken} and English \emph{milk} are Germanic cousins (Lewis \& Short). English \emph{emulsion} comes from the compound \emph{ēmulgēre}.
\end{cousinweb}

\begin{grammarlens}[title={What you've built}]
One more word for telling, explaining and describing at B1.
\end{grammarlens}

\subsection*{Your turn}
\begin{itemize}
\raggedright
  \item \emph{Say it:} \emph{mulgeō, mulgēre}
  \item \emph{Say it:} \emph{mulgeō, mulgēre}, once more
  \item \emph{Say it:} a sentence of your own with \emph{mulgeō, mulgēre}
  \item \emph{Recall:} say the Latin for a sheepfold, then the Latin for to shear, then say \emph{mulgeō, mulgēre} again
\end{itemize}

\begin{tcolorbox}[breakable,colback=teal!4,colframe=teal!35!black,title={Before you move on}]
What does mulgeō, mulgēre mean? (\textquotedblleft{}To milk\textquotedblright{}.) Say it once more.
\end{tcolorbox}
\section[(dialogue)]{Words from chapters 471–478 — a review}

\label{lesson:LA-R478-review-471-478}



\begin{quote}
Say the Latin for a lathe and to turn on a lathe.
\end{quote}

\subsection*{Your turn}
Say these aloud:

\begin{itemize}
\raggedright
  \item a lathe, to turn on a lathe, to polish, to sharpen, a stonecutter
  \item a distaff, a spindle, a web on the loom, to spin, wool-working
  \item the countryside, a farm, an ox stall, a cottage, a hut
  \item to plough, a furrow, a ploughshare, a ploughman, to harrow
  \item the harvest, a reaper, a sickle, an ear of grain, wheat
  \item to thresh, to winnow, chaff, a stalk, a granary
  \item livestock, a flock, a herd, to pasture, a pasture
  \item a shepherd, an ox-driver, a sheepfold, to shear, to milk
\end{itemize}

\begin{tcolorbox}[breakable,colback=teal!4,colframe=teal!35!black,title={Before you move on}]
A lathe? (\textbf{tornus, tornī}.) To turn on a lathe? (\textbf{tornō, tornāre}.) To polish? (\textbf{poliō, polīre}.) To sharpen? (\textbf{acuō, acuere}.) A stonecutter? (\textbf{lapicīda, lapicīdae}.) A distaff? (\textbf{colus, colī}.) A spindle? (\textbf{fūsus, fūsī}.) A web on the loom? (\textbf{tēla, tēlae}.) To spin? (\textbf{neō, nēre}.) Wool-working? (\textbf{lānificium, lānificiī}.) The countryside? (\textbf{rūs, rūris}.) A farm? (\textbf{fundus, fundī}.) An ox stall? (\textbf{būbīle, būbīlis}.) A cottage? (\textbf{casa, casae}.) A hut? (\textbf{tugurium, tuguriī}.) To plough? (\textbf{arō, arāre}.) A furrow? (\textbf{sulcus, sulcī}.) A ploughshare? (\textbf{vōmer, vōmeris}.) A ploughman? (\textbf{arātor, arātōris}.) To harrow? (\textbf{occō, occāre}.) The harvest? (\textbf{messis, messis}.) A reaper? (\textbf{messor, messōris}.) A sickle? (\textbf{falx, falcis}.) An ear of grain? (\textbf{spīca, spīcae}.) Wheat? (\textbf{trīticum, trīticī}.) To thresh? (\textbf{trītūrō, trītūrāre}.) To winnow? (\textbf{ventilō, ventilāre}.) Chaff? (\textbf{palea, paleae}.) A stalk? (\textbf{stipula, stipulae}.) A granary? (\textbf{horreum, horreī}.) Livestock? (\textbf{pecus, pecoris}.) A flock? (\textbf{grex, gregis}.) A herd? (\textbf{armentum, armentī}.) To pasture? (\textbf{pāscō, pāscere}.) A pasture? (\textbf{pāscuum, pāscuī}.) A shepherd? (\textbf{ōpiliō, ōpiliōnis}.) An ox-driver? (\textbf{bubulcus, bubulcī}.) A sheepfold? (\textbf{caulae, caulārum}.) To shear? (\textbf{tondeō, tondēre}.) To milk? (\textbf{mulgeō, mulgēre}.)
\end{tcolorbox}
